% ─────────────────────────────────────────────────────────────────────────
% §2 Related Work
% ─────────────────────────────────────────────────────────────────────────
\section{Related Work}\label{sec:related}

Three strands of existing literature bear on this paper: (i) RLHF-based
alignment and the engineering of refusal behavior, (ii) comparative
audits of LLM political and safety behavior, and (iii) studies of
LLM agent populations in computational social science. We survey each
briefly and position our contribution relative to them.

\subsection{RLHF Alignment and Refusal Engineering}\label{sec:related-rlhf}

Reinforcement Learning from Human Feedback \citep{christiano2017deep,
ouyang2022training,bai2022training} has become the dominant
post-training pipeline for commercial LLMs. A substantial body of work
documents that RLHF does not only make models more \emph{helpful}: it
also systematically shapes their refusal behavior, typically in ways
that trade off engagement against safety \citep{bai2022constitutional,
ganguli2022red,perez2022discovering}. The alignment literature
generally frames refusal as the intended output on a recognizable class
of unsafe prompts, but recent work has begun to characterize the
\emph{over-refusal} phenomenon --- models refusing benign prompts
superficially resembling unsafe ones --- as an empirical concern
\citep{rottger2024xstest,cui2024orbench}.

Our work contributes to this strand by exposing a richer taxonomy than
the binary \emph{refuse/comply} distinction typically used. We
distinguish \emph{hard refusal} (template blocking), \emph{soft refusal}
(substantive engagement that hedges, punts, or reframes), and \emph{on-task}
responses, and find that soft refusal dominates the refusal category
across all five vendors. This is consistent with the intuition that
contemporary RLHF optimizes for engaged hedging rather than for
recognizable blocking, and suggests that binary over-refusal metrics
may understate the prevalence of vendor-level alignment shaping.

\subsection{Comparative Audits of LLM Behavior}\label{sec:related-audits}

A growing set of papers compares LLM behavior across vendors on
specific domains. \citet{santurkar2023whose} and \citet{durmus2024opinionqa}
probe U.S.\ political opinion alignment across vendor families,
reporting systematic vendor-dependent biases on polarized policy
questions. \citet{rottger2024safetyprompts} benchmarks refusal on a
multilingual safety prompt set, with coverage across English-speaking
LLM vendors. \citet{hartmann2023political} documents ideological
leaning differences in ChatGPT's responses, finding consistent
left-leaning bias under their prompt regime.

Chinese-vendor safety and political behavior has been studied more
recently. \citet{huang2023ceval} established general Chinese-language
capability benchmarks; more directly relevant,
\citet{naseh2025r1dacted} investigate local censorship in
\vendor{DeepSeek}'s R1 model, finding that R1 systematically refuses
politically sensitive prompts in ways not observed in Western models,
and characterize this as the first instance of local censorship
alignment at the model layer. Most closely related to the present paper,
\citet{ko2026bilingual} benchmarks 17 LLMs on Taiwan sovereignty
questions in parallel Chinese/English query pairs and finds significant
language-dependent bias in 15 of 17 models, with models behaving
differently on the same sovereignty question depending on query language.
These neighboring works establish Chinese-vendor and Taiwan-political
LLM behavior as a live empirical concern but do not report the layered
architectural decomposition or the topic-stratified vendor taxonomy that
this paper develops.

Our paper contributes to this strand a specifically Taiwan-political
audit, with explicit paired-prompt design (the same 200 prompts put
to each of five vendors) and statistical inference at a level not
common in existing audits. Most cited audits report point-estimate
divergences without confidence intervals; we report paired bootstrap
BCa CIs throughout and argue several findings from CI-disjoint evidence.
We also contribute the topic$\times$vendor$\times$layer stratification
(§\ref{sec:discussion-three-way}) as a reusable methodological frame.

\subsection{LLM Agent Simulation}\label{sec:related-simulation}

The use of LLMs as simulants for human populations in social science
has grown rapidly \citep{park2023generative,horton2023llm,
argyle2023samples,tornberg2023simulating}. The methodological
assumption underlying many such studies is that the simulated population's
responses reflect some recoverable underlying distribution --- e.g., of
political preferences, of deliberative stances --- that is faithful to
the target human population modulo appropriate prompt engineering.
Several authors have raised concern that this assumption is not robust
to vendor-level alignment shaping \citep{santurkar2023whose,
sun2024personasteered}: different vendors may produce systematically different
simulated populations from identical prompts.

Our paper provides direct empirical support for this concern in a
specific and well-defined setting. The four-profile vendor taxonomy
(§\ref{sec:results-finding5}) implies that an LLM agent simulation of
Taiwanese political engagement would produce materially different
sovereignty-topic behavior under different vendor choices --- with
\vendor{DeepSeek}-based populations dropping to 10\% sovereignty-topic
engagement versus \vendor{Grok}-based populations sustaining 82\%.
We discuss implications for simulation practice explicitly in
§\ref{sec:discussion-simulation} and frame the audit as a tool for
quantifying vendor-dependent confounds in downstream simulation studies.

\paragraph{Positioning.}
Relative to the cited literature, this paper distinguishes itself in
four ways. First, the prompt set is specifically designed around
Taiwan political sensitivity in Traditional Chinese, a domain with
limited public audit data. Second, we include \vendor{Kimi} and
\vendor{DeepSeek} in the vendor panel, covering Chinese-owned
vendors at their recommended production endpoints.
Third, the statistical protocol uses prompt-level paired bootstrap for
all primary quantities, and several key findings are CI-disjoint at
95\%. Fourth, the four-category response taxonomy (with API-blocked as a
distinct fourth category) is new to the LLM audit literature to our
knowledge and is crucial for surfacing the two-layer architecture
finding (§\ref{sec:results-finding4}).
